\documentclass[10pt,a4paper]{amsart}
\usepackage[margin=19mm]{geometry}
\usepackage{amsmath,amssymb,mathtools}
\usepackage[scr=rsfs,cal=euler]{mathalfa}
\usepackage{xfrac}
\usepackage[unicode,hidelinks,bookmarks=false]{hyperref}
\newcommand{\ie}{i.e.\ }
\newcommand{\cf}{cf.\ }
\newcommand{\resp}{respectively\ }
\newcommand{\Subsection}{\subsection}
\providecommand{\nobreakdash}{\nobreak\hbox{-}}
\setlength{\emergencystretch}{2em}
\multlinegap0pt
\usepackage{amssymb}
\usepackage{xcolor}
\newtheorem{thm}{Theorem}
\title{\textbf{Explicit Factorization of $x^{p+1}-1$ over $\mathbb{Z}_{p^e}$: \\ A Structural Approach via Dickson Polynomials}}
\author{Yongchao Wang, Yang Ding, Jiansheng Yang, Zhiqiu Huang}
\date{Source: arXiv:2604.19038v3 (CC BY 4.0)}
\begin{document}
\maketitle

\noindent\emph{Source and licence.} \textbf{Explicit Factorization of $x^{p+1}-1$ over $\mathbb{Z}_{p^e}$: \\ A Structural Approach via Dickson Polynomials}. Version arXiv:2604.19038v3 (CC BY 4.0), distributed by arXiv under the Creative Commons Attribution 4.0 International licence, http://creativecommons.org/licenses/by/4.0/, as recorded in the arXiv licence metadata for this exact version and shown on the abstract page https://arxiv.org/abs/2604.19038v3. Full licence notice: \texttt{licenses/arxiv-2604.19038-CC-BY-4.0.txt}.\par\smallskip

\noindent\textit{Editorial scope.} Counted unit: the article's thm thm:base\_factorization (source0.tex lines 201--219). The article's complete original proof of it is delivered with it (source0.tex lines 222--238, one \texttt{proof} environment of 17 lines). No mathematics is written, repaired or re-derived: the statement and the proof are the article's own lines, in the article's own notation, and no retained line is substituted or re-flowed.  No further source-local context is needed: every cross-reference of every retained line resolves inside the two delivered spans.  Nothing was omitted from any retained span, so the package carries no omission record.  No retained line contains a citation, so the wrapper carries no bibliography and no bibliography entry.  No retained line contains a figure environment, an image-inclusion command or drawing code, so no image is delivered and the package declares no asset dependency.  Editorial formatting only, in addition: an AMS \texttt{amsart} preamble that restates the article's own declarations and macros used by the retained lines, namely the environments \texttt{thm} on separate counters, together with the standard packages those lines need.  The retained environments print in this document as Theorem 1 in order of appearance, whereas the article prints the same items with the numbering of its own full document.  The complete native source of the article as distributed in the arXiv e-print for this exact version is bundled unchanged as \texttt{sources/198840/source0.tex}.
\begin{thm}\label{thm:base_factorization}
The irreducible factorization of $x^{p+1}-1$ over $\mathbb{F}_p$ is given by:
\begin{equation}
x^{p+1}-1 = (x-1)(x+1) \Psi(x) \prod_{i \in S'} (x^2 - A_i x + 1),
\end{equation}
where:
\begin{itemize}
    \item The linear factors $(x-1)$ and $(x+1)$ correspond to the cosets $\{0\}$ and $\{(p+1)/2\}$.
    \item The term $\Psi(x)$ handles the case of $x^2+1$:
    $$ \Psi(x) = \begin{cases} x^2+1 & \text{if } p \equiv 3 \pmod 4 \\ \text{1 (empty product)} & \text{if } p \equiv 1 \pmod 4 \end{cases} $$
    \item $S'$ is the set of representatives for the remaining cosets of size 2, defined as $S' = \{1, 2, \dots, \frac{p+1}{2}\} \setminus \{\frac{p+1}{4}\}$ if $p \equiv 3 \pmod 4$, and $S' = \{1, 2, \dots, \frac{p+1}{2}\}$ if $p \equiv 1 \pmod 4$. Note that the index $\frac{p+1}{4}$ corresponds to the factor $x^2+1$ (where $A_i=0$). The coefficients $A_i$ for these factors are explicitly generated as follows:
    \begin{enumerate}
        \item \textbf{Seed Generation:} Let $f(x) = x^2 - a_1 x + a_2 \in \mathbb{F}_p[x]$ be a primitive polynomial (where $a_1$ and $a_2$ correspond to the trace and norm of the primitive element $\beta \in \mathbb{F}_{p^2}$, respectively). The generator $A_1$ is obtained via:
        $$A_1 = D_{p-1}(a_1, a_2).$$
        \item \textbf{Recursive Expansion:} For $i \ge 2$, coefficients are generated by the recurrence:
        $$A_i = D_i(A_1, 1) = A_1 A_{i-1} - A_{i-2}.$$
    \end{enumerate}
\end{itemize}
\end{thm}

\begin{proof}
The factorization structure in (i) and the specific linear/quadratic forms follow directly from the size of the $p$-cyclotomic cosets modulo $p+1$, which is at most 2 since $p^2 \equiv 1 \pmod{p+1}$.

We focus on proving the generative mechanism for $A_i$. Let $\beta$ be a root of the primitive polynomial $f(x) = x^2 - a_1 x + a_2$ in $\mathbb{F}_{p^2}$. Since $f(x)$ is primitive, $\beta$ is a generator of $\mathbb{F}_{p^2}^*$. We set $\alpha = \beta^{p-1}$. Since the order of $\beta$ is $p^2-1$, the order of $\alpha$ is $(p^2-1)/(p-1) = p+1$, making $\alpha$ a primitive $(p+1)$-th root of unity.

The coefficient for the $i$-th factor is $A_i = \alpha^i + \alpha^{-i}$. Recall the fundamental property of Dickson polynomials (Waring's formula): $u^k + v^k = D_k(u+v, uv)$.
\begin{itemize}
    \item \textbf{For the seed $A_1$:} We have $A_1 = \alpha + \alpha^{-1} = \beta^{p-1} + (\beta^{-1})^{p-1}$. Note that in $\mathbb{F}_{p^2}$, the conjugate of $\beta$ is $\beta^p$. Thus, the term $\alpha^{-1}$ corresponds to the conjugate action on the primitive root relative to the norm. Specifically, we observe:
    $$A_1 = \beta^{p-1} + (\beta^p)^{p-1} = D_{p-1}(\beta+\beta^p, \beta^{p+1}).$$
    From Vieta's formulas on $f(x)$, we have $\beta + \beta^p = a_1$ and $\beta^{p+1} = a_2$. Therefore, $A_1 = D_{p-1}(a_1, a_2)$.
    
    \item \textbf{For the recursion:} Since $A_i = \alpha^i + (\alpha^{-1})^i$ and $\alpha \cdot \alpha^{-1} = 1$, we simply apply the Dickson property again with $u=\alpha, v=\alpha^{-1}$:
    $$A_i = D_i(\alpha + \alpha^{-1}, \alpha \alpha^{-1}) = D_i(A_1, 1).$$
    This satisfies the recurrence $D_i(x, 1) = x D_{i-1}(x, 1) - D_{i-2}(x, 1)$.
\end{itemize}
This confirms that the entire set of coefficients is deterministically generated by the Dickson mapping of the primitive polynomial's coefficients.
\end{proof}

\end{document}
